% Generated from validated Article Draft Result. Do not hand-edit; revise the structured draft instead.
\noindent\textbf{7月のvLLM、SGLang、FlashInferは、model supportを増やすだけでなく、default execution path、speculative decoding、MoE・Blackwell向けkernelといった実行系の中核を更新した。研究側でもdiffusion servingとMoE-aware speculative decodingが、次のbottleneckを別の角度から捉えている。}\autocite{src-50e29a09d75a,src-5a9041d41ad0,src-76c7b104c7df,src-b80517d03022}

\subsection{serving frameworkの更新は「対応モデル追加」だけではない}

vLLM v0.25.0はdense model向けModel Runner V2をdefault execution pathへ切り替え、legacy PagedAttention implementationを削除した。続くv0.26.0では新しいmodel family supportに加え、attention、speculative decoding、KV offload、DeepSeek-V4関連のoptimizationが拡張された。7月の変更は、互換性リストの更新と内部execution architectureの更新が同時に進む例である。\autocite{src-76c7b104c7df,src-b811cc97eff4}


\begin{claimboundary}
vLLM release noteのperformance percentageは特定kernel・model条件でのproject measurementである。またintegration supportは、そのmodel自体のqualityやvLLM外でのavailabilityを証明しない。\autocite{src-76c7b104c7df,src-b811cc97eff4}
\end{claimboundary}

SGLang v0.5.15はSpec V2 pathをdefault化し、新しいmodel/hardware combination向けのproduction tuningを含めた。v0.5.16ではdraft confidenceに応じてverification windowを変えるDSparkを導入した。speculative decodingが単なるoptionではなく、default pathとalgorithm designの双方で更新されている点が重要である。\autocite{src-5a9041d41ad0,src-bf2a6d1b9b56}


\begin{claimboundary}
SGLangが示すthroughputやTPSはbatch size、tensor parallel、model、GPU構成に依存するproject benchmarkである。本誌ではalgorithm/releaseの存在と、一般化されたspeed superiorityを分離する。\autocite{src-5a9041d41ad0,src-bf2a6d1b9b56}
\end{claimboundary}

FlashInfer v0.6.15はexpert-parallel MoE supportをdefault installへ入れ、Blackwell familyのserving coverageを拡張した。v0.6.16ではMegaMoE kernels、さらに広いBlackwell model coverage、unified MoE API、distributed servingの変更を追加した。frontier model側でsparse MoEが増えるほど、kernelとcommunicationの設計がserving stackの前景へ出てくる。\autocite{src-020bbc9ce508,src-b80517d03022}


\begin{claimboundary}
FlashInferのhardware/model supportはGPU generationへの依存が大きく、release note中のthroughputは特定条件のproject benchmarkである。Blackwell対応の拡大を、そのまま任意環境の性能保証とは扱わない。\autocite{src-020bbc9ce508,src-b80517d03022}
\end{claimboundary}

\subsection{researchは次のcost modelを探っている}

研究側ではFlashDiffがdiffusion-model servingでregional executionとschedulingを提案し、毎回full regionを計算・通信するcostを減らす方向を示す。一方EcoSpecはsparse MoEのspeculative decodingでtoken acceptanceだけでなくexpert activation/scatteringのcostを考慮する。どちらも『何token通ったか』だけでは捉えにくい、model構造固有のexecution costをschedulerやspeculationへ持ち込む試みである。\autocite{src-50e29a09d75a,src-d6756f2b6aa6}


\begin{claimboundary}
FlashDiffとEcoSpecはいずれもauthor-evaluated systems researchであり、報告されたgainは選択model、hardware、workload、baselineに依存する。ここではproduction superiorityではなく、serving設計が扱おうとしている新しいcost要因のsignalとして用いる。\autocite{src-50e29a09d75a,src-d6756f2b6aa6}
\end{claimboundary}

\subsection{model releaseを支える「外側」の競争}

vLLMのrunner、SGLangのspeculation、FlashInferのMoE kernelsは異なるlayerを触っているが、共通してmodel architectureとhardwareの変化をruntimeへ吸収する仕事である。7月のfrontier-model更新を実際のserviceへ載せるには、model weightやAPIだけでなく、attention、KV、expert routing、kernel、schedulerまで追う必要がある。月次SpecialでInference \& Servingを独立テーマにする理由はここにある。\autocite{src-020bbc9ce508,src-b811cc97eff4,src-bf2a6d1b9b56,src-d6756f2b6aa6}
